\batchmode
\documentclass{article}
\makeatletter

\oddsidemargin=0in
\textwidth=6.5in
\topmargin=-0.5in
\textheight=9in

\parindent=0in

\begin{document}

 {\bf Mathematical Modeling (Math512 sec. 010 \& 080) \hfill Fall 2001 \\
MW 1430-1520 Mem 108 \& F 1430-1520 Mem 109} \\
\copyright 2001 L. F. Rossi \hspace{0.25in} All rights reserved.

\bigskip

 Prof. L. F. Rossi \\
Office: Ewg 524 \\
Telephone: x1880 \\
Email: {\tt rossi@math.udel.edu} \\
WWW: {\tt http://www.math.udel.edu/$\sim$ rossi}

\bigskip

 {\bf Your objective:} You will learn to use calculus, linear algebra,
differential equations and everything else you have learned to solve
everyday problems.  Ultimately, you must develop an intuition and
resourcefulness about solving problems.

\bigskip

 {\bf Your resources:} All of the following will help you achieve your
objectives:
\begin{list}{$\bullet$ }{}

\item Time: Attendance in this course is mandatory.

\item Office hours: Mon 1300-1400, Wed 1300-1400, Thurs 1000-1100 
or by appointment.
Office hours are one of the most valuable and least used resources at
the University, and I hope you will take advantage of them.  I want to
help you learn this material, so do not be shy about seeing me outside
of class.  Your exams will only be distributed during office hours.
If you need to see me at a time other than an office hour, feel free
to ``drop in'' or make an appointment.  

\item Your classmates: Math is not a competitive sport.  There are
many obvious reasons to work together.  Even if you end up helping
others most of the time, teaching is one of the best ways to gain a
deeper understanding of a subject.  In fact, half of your grade in
this course is a team effort.

\item Textbooks \& library research: 
There is no required textbook for this class, but you will want to
keep a few books handy.  First, if your fundamentals are a bit rusty,
you should dust off your calculus textbook.  Next, a standard text on
ordinary differential equations and linear algebra will definitely be
useful.  I will play general modeling texts on reserve in the
library.  Also, I
expect students to visit the library frequently to fill in gaps in
their knowledge.

\end{list}

\bigskip

 {\bf Themes:}

 This course will have several common themes that run through every
subject.

\begin{itemize}

\item Modeling: What is a model?  What is not?  
Is the model under consideration sufficient to answer the question at
hand?  How can we assess the strengths and weaknesses of the model?

\item What is the structure of an effective solution to a problem?
Who is the audience?  What is your objective and what role does the
audience play?

\end{itemize}

 {\bf Tentative topics:}

\begin{itemize}

\item Scaling, dimensional analysis and dimensionless constants:
Converting units, scale experiments, dynamic similarity...

\item Working with data I.  Interpolation, least squares, parameter
estimation, splines.

\item Conservations laws: Partial differential equations, 
flux, diffusion, wave propagation, 
Helmholtz equation, traffic, pipe flows...

\item Working with data II.  Assessing the quality of a model.  Norms and 
metrics.

\item Coding theory.  The football pool problem.  The colored hat
problem. Hamming codes.

\item Signal processing, frequency domain, scale domain, sampling,
bias, etc.

\end{itemize}

\newpage

 {\bf Regular problem sets:}

 At regular intervals of approximately one week, you will be required
to complete problem sets.

\bigskip

 {\bf Projects: }

 You will be required to complete three projects.  Projects 
are open-ended questions taken from previous Mathematical Contest in
Modeling problems and other sources.  There is no ``answer in the back 
of the book.''  You will be permitted to work in groups of up to three 
people on these problems, and there will be individual oral exams
based on your projects.

 The first two projects will have the form of a full written report.
This report should have an complete introduction with all relevant
background information, a statement of assumptions, an analysis and
solution the problem, an assessment of the quality of your solution
and a conclusion.  The last project presentation will be oral followed
by an on-site comparison of competing solutions.
I expect all projects to be complete and concise.

\bigskip

 {\bf Final exam: } 

 There will be a brief final exam at the end of the
semester on traditional topics surveyed in the course.

\bigskip

 {\bf Grading policy:}  

 Your grade is determined solely by your
participation, problem sets and final project. \\
{\center \begin{tabular}{ll}
Homework & 20\% \\
Class participation & 10\% \\
Projects & 30\% \\
Oral exam & 30\% \\
Final exam & 10\%
\end{tabular} \\}
Final letter grades will be assigned strictly based on the following
percentages of your total point score: \\
\documentclass{article}

\oddsidemargin=0in
\textwidth=6.5in
\topmargin=-0.5in
\textheight=9in

\parindent=0in
\par\usepackage[dvips]{color}
\pagecolor[gray]{.7}


\makeatletter
\count@=\the\catcode`\_ \catcode`\_=8 
\newenvironment{tex2html_wrap}{}{} \catcode`\_=\count@
\makeatother
\let\mathon=$
\let\mathoff=$
\ifx\AtBeginDocument\undefined \newcommand{\AtBeginDocument}[1]{}\fi
\newbox\sizebox
\setlength{\hoffset}{0pt}\setlength{\voffset}{0pt}
\addtolength{\textheight}{\footskip}\setlength{\footskip}{0pt}
\addtolength{\textheight}{\topmargin}\setlength{\topmargin}{0pt}
\addtolength{\textheight}{\headheight}\setlength{\headheight}{0pt}
\addtolength{\textheight}{\headsep}\setlength{\headsep}{0pt}
\setlength{\textwidth}{349pt}
\newwrite\lthtmlwrite
\makeatletter
\let\realnormalsize=\normalsize
\global\topskip=2sp
\def\preveqno{}\let\real@float=\@float \let\realend@float=\end@float
\def\@float{\let\@savefreelist\@freelist\real@float}
\def\end@float{\realend@float\global\let\@freelist\@savefreelist}
\let\real@dbflt=\@dbflt \let\end@dblfloat=\end@float
\let\@largefloatcheck=\relax
\def\@dbflt{\let\@savefreelist\@freelist\real@dbflt}
\def\adjustnormalsize{\def\normalsize{\mathsurround=0pt \realnormalsize
 \parindent=0pt\abovedisplayskip=0pt\belowdisplayskip=0pt}\normalsize}%
\def\lthtmltypeout#1{{\let\protect\string\immediate\write\lthtmlwrite{#1}}}%
\newcommand\lthtmlhboxmathA{\adjustnormalsize\setbox\sizebox=\hbox\bgroup}%
\newcommand\lthtmlvboxmathA{\adjustnormalsize\setbox\sizebox=\vbox\bgroup%
 \let\ifinner=\iffalse }%
\newcommand\lthtmlboxmathZ{\@next\next\@currlist{}{\def\next{\voidb@x}}%
 \expandafter\box\next\egroup}%
\newcommand\lthtmlmathtype[1]{\def\lthtmlmathenv{#1}}%
\newcommand\lthtmllogmath{\lthtmltypeout{l2hSize %
:\lthtmlmathenv:\the\ht\sizebox::\the\dp\sizebox::\the\wd\sizebox.\preveqno}}%
\newcommand\lthtmlfigureA[1]{\let\@savefreelist\@freelist
       \lthtmlmathtype{#1}\lthtmlvboxmathA}%
\newcommand\lthtmlfigureZ{\lthtmlboxmathZ\lthtmllogmath\copy\sizebox
       \global\let\@freelist\@savefreelist}%
\newcommand\lthtmldisplayA[1]{\lthtmlmathtype{#1}\lthtmlvboxmathA}%
\newcommand\lthtmldisplayB[1]{\edef\preveqno{(\theequation)}%
  \lthtmldisplayA{#1}\let\@eqnnum\relax}%
\newcommand\lthtmldisplayZ{\lthtmlboxmathZ\lthtmllogmath\lthtmlsetmath}%
\newcommand\lthtmlinlinemathA[1]{\lthtmlmathtype{#1}\lthtmlhboxmathA  \vrule height1.5ex width0pt }%
\newcommand\lthtmlinlineA[1]{\lthtmlmathtype{#1}\lthtmlhboxmathA}%
\newcommand\lthtmlinlineZ{\egroup\expandafter\ifdim\dp\sizebox>0pt %
  \expandafter\centerinlinemath\fi\lthtmllogmath\lthtmlsetinline}
\newcommand\lthtmlinlinemathZ{\egroup\expandafter\ifdim\dp\sizebox>0pt %
  \expandafter\centerinlinemath\fi\lthtmllogmath\lthtmlsetmath}
\def\lthtmlsetinline{\hbox{\vrule width.1em\vtop{\vbox{%
  \kern.1em\copy\sizebox}\ifdim\dp\sizebox>0pt\kern.1em\else\kern.3pt\fi
  \ifdim\hsize>\wd\sizebox \hrule depth1pt\fi}}}
\def\lthtmlsetmath{\hbox{\vrule width.1em\vtop{\vbox{%
  \kern.1em\kern0.8 pt\hbox{\hglue.17em\copy\sizebox\hglue0.8 pt}}\kern.3pt%
  \ifdim\dp\sizebox>0pt\kern.1em\fi \kern0.8 pt%
  \ifdim\hsize>\wd\sizebox \hrule depth1pt\fi}}}
\def\centerinlinemath{%\dimen1=\ht\sizebox
  \dimen1=\ifdim\ht\sizebox<\dp\sizebox \dp\sizebox\else\ht\sizebox\fi
  \advance\dimen1by.5pt \vrule width0pt height\dimen1 depth\dimen1 
 \dp\sizebox=\dimen1\ht\sizebox=\dimen1\relax}

\def\lthtmlcheckvsize{\ifdim\ht\sizebox<\vsize\expandafter\vfill
  \else\expandafter\vss\fi}%
\makeatletter \tracingstats = 1 


\begin{document}
\pagestyle{empty}\thispagestyle{empty}%
\lthtmltypeout{latex2htmlLength hsize=\the\hsize}%
\lthtmltypeout{latex2htmlLength vsize=\the\vsize}%
\lthtmltypeout{latex2htmlLength hoffset=\the\hoffset}%
\lthtmltypeout{latex2htmlLength voffset=\the\voffset}%
\lthtmltypeout{latex2htmlLength topmargin=\the\topmargin}%
\lthtmltypeout{latex2htmlLength topskip=\the\topskip}%
\lthtmltypeout{latex2htmlLength headheight=\the\headheight}%
\lthtmltypeout{latex2htmlLength headsep=\the\headsep}%
\lthtmltypeout{latex2htmlLength parskip=\the\parskip}%
\lthtmltypeout{latex2htmlLength oddsidemargin=\the\oddsidemargin}%
\makeatletter
\if@twoside\lthtmltypeout{latex2htmlLength evensidemargin=\the\evensidemargin}%
\else\lthtmltypeout{latex2htmlLength evensidemargin=\the\oddsidemargin}\fi%
\makeatother
{\newpage\clearpage
\lthtmlinlinemathA{tex2html_wrap_inline30}%
$\sim$%
\lthtmlinlinemathZ
\hfill\lthtmlcheckvsize\clearpage}

{\newpage\clearpage
\lthtmlinlinemathA{tex2html_wrap_inline32}%
$\bullet$%
\lthtmlinlinemathZ
\hfill\lthtmlcheckvsize\clearpage}

{\newpage\clearpage
\lthtmlinlinemathA{tex2html_wrap_inline14}%
$\mbox{$\bullet$}$%
\lthtmlinlinemathZ
\hfill\lthtmlcheckvsize\clearpage}

{\newpage\clearpage
\lthtmlinlinemathA{tex2html_wrap_inline38}%
$\leftarrow$%
\lthtmlinlinemathZ
\hfill\lthtmlcheckvsize\clearpage}

{\newpage\clearpage
\lthtmlinlinemathA{tex2html_wrap_inline40}%
$\rightarrow$%
\lthtmlinlinemathZ
\hfill\lthtmlcheckvsize\clearpage}


\end{document}
